% Laget automatisk av elevark/lag_elevark.py fra opplegg/video-muffins.html. Endringer her blir overskrevet.
\documentclass[11pt]{article}
\usepackage{elevark}
\begin{document}
\elevarkhode{Fysikk 1 · Bevegelse og krefter}{Videoanalyse av muffinsformer}{https://fysikklab.no/opplegg/video-muffins.html}{Opplegget på nett}

En muffinsform når fort terminalfarten, der luftmotstanden er like stor som tyngden. Når dere legger flere former i hverandre, blir massen større, men formen den samme. Da kan dere teste modellen $L = kv^{2}$ og bestemme luftmotstandstallet $k$.

\section*{Hensikt}

\begin{itemize}
\item Finne terminalfarten ved videoanalyse.
\item Undersøke om luftmotstanden er proporsjonal med $v^{2}$.
\item Bestemme luftmotstandstallet $k$.
\end{itemize}

\section*{Utstyr}

\begin{itemize}
\item Flere like muffinsformer av papir, og en vekt som måler med minst 0,1~g nøyaktighet.
\item Mobil eller kamera på stativ, en målestokk som henger loddrett, og Tracker.
\item Slipp fra omtrent 2~m høyde, for eksempel fra en stol eller trapp.
\end{itemize}

\section*{Framgangsmåte}

\begin{enumerate}
\item Vei én muffinsform. Vei gjerne ti former og del på ti for å få mer nøyaktig masse.
\item Film én form som slippes fra ro foran målestokken. Hold formen med åpningen opp, og slipp den uten å dytte.
\item Gjenta med 2, 3 og 4 former stablet i hverandre. Formen og størrelsen er den samme, men massen øker.
\item I Tracker: Kalibrer, sett aksene med $y$ oppover, og spor formen. Kopier $t$ og $y$ inn i verktøyet på nettsiden.
\item Finn tidsrommet der grafen er en rett linje, og tilpass en rett linje der. Stigningstallet er terminalfarten. Sett den inn i tabellen.
\end{enumerate}

Uten egen video: Tracker-eksemplet «CupsClips» viser muffinsformer med ulikt antall binders. Ifølge beskrivelsen i fila veier formen med to binders omtrent 1,4~g og formen med seks binders omtrent 3,1~g. Dere kan også øve med simuleringen Videoanalyse uten video\footnote{\url{fysikklab.no/sim/videoanalyse.html}}.

\section*{Teorien}

Ved terminalfarten er $L = G$, altså $kv_{\mathrm{t}}^{2} = mg$. Da er $k = mg / v_{\mathrm{t}}^{2}$. Hvis modellen stemmer, er $v_{\mathrm{t}}^{2}$ proporsjonal med $m$, og grafen av $v_{\mathrm{t}}^{2}$ mot $m$ blir en rett linje gjennom origo med stigningstall $g/k$.

\section*{Oppgaver}

\begin{oppgaveliste}
\oppgave \fase{gjett}{Gjett}Last inn eksempeldataene for én form. Hvor lang tid tar det før formen har nådd terminalfarten? Hvordan ser du det?
\svar
\oppgave \fase{test}{Mål}Finn terminalfarten for 1, 2, 3 og 4 former. Velg tidsrom fra omtrent 0,6~s og ut. Sett dem inn i tabellen, med massen 0,40~g per form.
\svar
\oppgave \fase{forklar}{Forklar}Er $k$ omtrent den samme for alle radene? Ligger punktene på en rett linje gjennom origo? Hva sier det om modellen $L = kv^{2}$?
\svar
\oppgave \fase{forklar}{Forklar}Hva ville grafen av $v_{\mathrm{t}}$ mot $m$ vist hvis luftmotstanden i stedet var proporsjonal med $v$?
\svar
\oppgave \fase{test}{Regn}Luftmotstandstallet kan skrives $k = \tfrac{1}{2}ρC_{\mathrm{d}}A$. Mål diameteren til formen, og regn ut formfaktoren $C_{\mathrm{d}}$. Bruk $ρ = 1{,}2\,\mathrm{kg/m^{3}}$.
\svar
\oppgave \fase{utfordring}{Utfordring}Hvorfor er det viktig at formene er stablet i hverandre, og ikke for eksempel tapet sammen side om side?
\svar
\end{oppgaveliste}

\end{document}
